\documentclass[12pt,a4paper]{article}
\usepackage{geometry}
\usepackage{titling}
\usepackage{fontspec}
\usepackage{unicode-math}
\usepackage{amsmath}
\usepackage{pgfplots}
\pgfplotsset{compat=1.16}
\usepackage{booktabs}
\usepackage{setspace}
\usepackage{relsize}
\setlength{\parindent}{0pt}
\setlength{\parskip}{0.7em}

\geometry{a4paper, top=2.2cm, bottom=3cm, left=2.5cm, right=2.5cm}
\setlength{\droptitle}{-3cm}

\definecolor{charcoal}{HTML}{36454F}

\iftrue
%\setstretch{1.025} % a little more space between lines
\setmainfont{texgyrepagella}[
  Extension = .otf,
  UprightFont = *-regular,
  BoldFont = *-bold,
  ItalicFont = *-italic,
  BoldItalicFont = *-bolditalic,
  Color=charcoal,
]
\fi

\setmathfont{texgyrepagella-math.otf}[Color=charcoal]
\pagecolor{white!6}

\pretitle{\begin{center}\Large\addfontfeatures{LetterSpace=-1.5}}
\posttitle{\end{center}\vspace*{-0.5em}}
\preauthor{\begin{center}\large\addfontfeatures{LetterSpace=-1.5}}
\postauthor{\end{center}\vspace*{-0.6cm}}

\title{A Short Note on Surprisal, Entropy, and Cross-Entropy Loss\\\vspace*{2pt}[Draft 1]}
\author{Devi Prasad M\\Manipal School of Information Sciences}
\date{October 2026}

\begin{document}
\maketitle

\section{Surprisal of an outcome}

Let $\symbf{q}$ be a discrete probability distribution over
possible outcomes, with $q_k = \symbf{q}(X = k)$.
The \emph{surprisal} of outcome $k$ under $\symbf{q}$ is
\begin{equation}
  h_{\symbf{q}}(k) = \log \dfrac{1}{q_k} = -\log q_k.
  \label{eq:surprisal}
\end{equation}
If $q_k = 0$, we define $h_{\symbf{q}}(k) = +\infty$.
This is a quantity for \emph{one specified outcome}.
We will see later that entropy is the average surprisal
over outcomes.

Throughout this note, $\log$ means the natural logarithm
$\ln$, so surprisal is measured in nats.

Why the minus sign? Let us first look at $\log q_k$
on its own. Figure~\ref{fig:ln} shows it for $q_k \in (0, 1]$.

\begin{figure}[ht]
  \centering
  \begin{tikzpicture}
    \begin{axis}[
      width=0.7\textwidth, height=5.5cm,
      xlabel={$q_k$}, ylabel={$\log q_k$},
      xmin=0, xmax=1, ymin=-5, ymax=0.8,
      domain=0.01:1, samples=100, grid=major,
      axis x line=middle, axis y line=left,
      xticklabel style={anchor=south, yshift=2pt},
      xlabel style={at={(ticklabel cs:1.03)}, anchor=west}]
      \addplot[thick] {ln(x)};
    \end{axis}
  \end{tikzpicture}
  \caption{The curve $\log q_k$ for $q_k \in (0, 1]$.
  It is $0$ at $q_k = 1$, is negative everywhere
  else, and drops without bound as $q_k \to 0^{+}$.}
  \label{fig:ln}
\end{figure}

Two things are clear from the curve.

\begin{itemize}
  \item \textbf{Direction:} $\log q_k$ goes
  \emph{up} as $q_k$ goes up. But surprise should
  go the other way: a likely outcome should be
  unsurprising, and a rare one very surprising.
  \item \textbf{Sign:} $\log q_k \le 0$. But we
  want surprise to be a non-negative number,
  like a length.
\end{itemize}

Putting a minus sign in front flips the curve
about the horizontal axis, and this fixes both
problems:
\begin{itemize}
  \item $q_k = 1$ gives surprisal $0$. A certain
        outcome is not surprising at all.
  \item $q_k \to 0^{+}$ gives surprisal
        $\to \infty$. A nearly impossible outcome
        is very surprising.
\end{itemize}

\newpage

Figure~\ref{fig:surprisal} shows the flipped curve.

\begin{figure}[ht]
  \centering
  \begin{tikzpicture}
    \begin{axis}[
      width=0.7\textwidth, height=5.5cm,
      xlabel={$q_k$}, ylabel={$-\log q_k$},
      xmin=0, xmax=1, ymin=-0.5, ymax=5,
      domain=0.01:1, samples=100, grid=major,
      axis x line=middle, axis y line=left,
      xlabel style={at={(ticklabel cs:1.03)}, anchor=west}]
      \addplot[thick] {-ln(x)};
    \end{axis}
  \end{tikzpicture}
  \caption{Surprisal $-\log q_k$ against $q_k$.
  It is $0$ at $q_k = 1$ and grows without bound
  as $q_k \to 0^{+}$.}
  \label{fig:surprisal}
\end{figure}

\section{Why surprisals add}

For two independent events $A$ and $B$,
probabilities multiply:
\begin{equation}
  P(A, B) = P(A)\,P(B).
  \label{eq:indep}
\end{equation}
We want a measure of surprise that \emph{adds}
for independent events, just like length or mass.
The logarithm turns a product into a sum.
Writing $h(A) = -\log P(A)$ for the surprisal
of an event, we get
\begin{equation}
  \begin{aligned}
    h(A, B) &= -\log P(A, B) = -\log\big(P(A)\,P(B)\big) \\
            &= -\log P(A) - \log P(B) = h(A) + h(B).
  \end{aligned}
  \label{eq:add}
\end{equation}
Thus, adding surprisals is nothing but the
multiplication rule of probability, transformed
via a logarithm. This is the same reason we
use log-likelihood for i.i.d. data, turning
a product of probabilities
into a sum.

\section{Entropy: average surprisal}

If we draw $N$ independent outcomes from $\symbf{q}$,
outcome $k$ comes up about $N q_k$ times, and each time
it brings surprisal $-\log q_k$. The average surprisal
per draw is therefore
\begin{equation}
  H(\symbf{q}) = \sum_{k} q_k \,(-\log q_k) = -\sum_{k} q_k \log q_k .
  \label{eq:entropy}
\end{equation}
The factor $q_k$ is simply the fraction of time outcome
$k$ actually happens. Every average needs such weights.

Here $\symbf{q}$ plays two roles: it gives the
\emph{weights}, and it also decides the \emph{surprisal}.

Outcomes with $q_k = 0$ contribute nothing. The term
$q_k(-\log q_k)$ is then of the form $0 \cdot \infty$,
which we take to be $0$ by convention
($0 \log 0 = 0$, justified because $q \log q \to 0$
as $q \to 0^{+}$).

\textbf{Example.} Take $\symbf{q} = (0.9,\, 0.1)$. Then
\begin{equation}
  H(\symbf{q}) = 0.9 \times 0.105 + 0.1 \times 2.303 \approx 0.325 \text{ nats}.
  \label{eq:example}
\end{equation}
The common outcome has small surprisal and large
weight, so it dominates the average.

\textbf{Two observations.}
\begin{enumerate}
    \item If one outcome has probability
            $1$, then $H = 0$.

    \item If all $K$ outcomes are equally
    likely, then $H = \log K$, the maximum.
    \begin{equation}
        H(\symbf{q})
        = \mathlarger{\sum_{k=1}^{K}} q_k \,(-\log q_k)
        = \mathlarger{\sum_{k=1}^{K}} \dfrac{1}{K} \bigg(-\log \dfrac{1}{K}\bigg)
        = K \cdot \dfrac{1}{K} \cdot \log K
        = \log K.
        \label{eq:uniform}
    \end{equation}
\end{enumerate}

\section{Cross-entropy: separating the two roles}

Now suppose outcomes really occur according to a
distribution $\symbf{p}$ (the truth), but we measure
surprise using our model $\symbf{q}$. Then the weights
come from $\symbf{p}$ and the surprisal comes
from $\symbf{q}$:
\begin{equation}
    H(\symbf{p}, \symbf{q}) = \sum_{k} p_k \,(-\log q_k) = -\sum_{k} p_k \log q_k .
    \label{eq:xent}
\end{equation}
\paragraph{In simple words:}
cross-entropy is the average surprise
of the model $\symbf{q}$ when the data actually comes
from $\symbf{p}$.

Replacing $\symbf{q} = \symbf{p}$ gives back
the entropy:
\begin{equation}
    H(\symbf{p}, \symbf{p}) = -\sum_{k} p_k \log p_k = H(\symbf{p}).
    \label{eq:xent-self}
\end{equation}

\section{One-hot labels}

In classification, the label is one-hot. If the true
class is $c$, then $p_c = 1$ and $p_k = 0$ for all
$k \neq c$. Every term with $p_k = 0$ vanishes,
and Eq.~\eqref{eq:xent} leaves us with
\begin{equation}
  H(\symbf{p}, \symbf{q}) = -\log q_c .
  \label{eq:onehot}
\end{equation}
So the cross-entropy loss for one example is simply
\emph{how surprised the model is by the correct class}.
If $q_c$ is near $1$, the loss is near $0$.
If $q_c$ is near $0$, the loss is large.

The probabilities given to wrong classes do not appear
directly. They matter only because $\symbf{q}$ sums
to $1$: whatever mass goes to wrong classes is
taken away from $q_c$.

\section{Surprisal, likelihood, and log-likelihood}

We often come across these five terms: surprisal,
likelihood, log-likelihood, entropy, and cross-entropy.
All these ideas are directly related to one another.
In this section we will connect them.

\paragraph{A data set of N examples.}
In training, we have $N$
examples $(x_1, c_1), \dots, (x_N, c_N)$,
where $c_i$ is the true class of input
$x_i$. Let $\theta$ denote the parameters
of the network. We write
$q_\theta(c_i \mid x_i)$ for the probability
that the network gives to the correct class of
example $i$. This is just the $q_c$ of
Eq.~\eqref{eq:onehot}, written for example $i$.

\paragraph{Likelihood.}
Assume the examples are independent. Then, by
Eq.~\eqref{eq:indep}, the probability that the network
gives to \emph{all} the observed labels together is a
product:
\begin{equation}
  P(c_1, \dots, c_N \mid x_1, \dots, x_N;\ \theta)
  = \prod_{i=1}^{N} q_\theta(c_i \mid x_i).
  \label{eq:lik}
\end{equation}
When we keep the data fixed and look at this number as a
function of $\theta$, we call it the \emph{likelihood} of
$\theta$. A good $\theta$ is one that makes the observed
labels probable, so we want the likelihood to be large.

\paragraph{Log-likelihood.}
The logarithm turns this product into a sum, just as in
Eq.~\eqref{eq:add}:
\begin{equation}
  \ell(\theta) = \log \prod_{i=1}^{N} q_\theta(c_i \mid x_i)
  = \sum_{i=1}^{N} \log q_\theta(c_i \mid x_i).
  \label{eq:loglik}
\end{equation}
This is the \emph{log-likelihood}. Each term is $\le 0$, so
$\ell(\theta) \le 0$. Since $\log$ is an increasing
function, making $\ell(\theta)$ large is the same as making
the likelihood large.

There is also a practical reason to prefer the log. A
product of thousands of numbers below $1$ becomes so small
that a computer rounds it to $0$. The sum in
Eq.~\eqref{eq:loglik} has no such problem.

\paragraph{Negative log-likelihood is total surprisal.}
Now put a minus sign in front of Eq.~\eqref{eq:loglik}.
Each term becomes a surprisal, Eq.~\eqref{eq:surprisal}:
\begin{equation}
  -\ell(\theta)
  = \sum_{i=1}^{N} \big(-\log q_\theta(c_i \mid x_i)\big)
  = \text{total surprisal of the observed labels}.
  \label{eq:nll}
\end{equation}
This quantity is called the \emph{negative log-likelihood}
(NLL). So surprisal and log-likelihood are the same numbers
with opposite signs: log-likelihood measures how
\emph{plausible} the model finds the data, and surprisal
measures how \emph{surprised} the model is by it.

\paragraph{Average surprisal is the cross-entropy loss.}
Divide Eq.~\eqref{eq:nll} by $N$:
\begin{equation}
  \bar{L}(\theta)
  = -\frac{1}{N}\,\ell(\theta)
  = \frac{1}{N}\sum_{i=1}^{N} \big(-\log q_\theta(c_i \mid x_i)\big).
  \label{eq:meanloss}
\end{equation}
Each term in the sum is the one-hot cross-entropy of
Eq.~\eqref{eq:onehot} for one example. So $\bar{L}(\theta)$
is the average cross-entropy loss over the data set, which
is exactly the number that deep learning libraries report.

Putting it all together, these four statements are the
same:
\begin{center}
\begin{tabular}{rl}
         & maximise the likelihood \\[2pt]
$\iff$   & maximise $\ell(\theta)$ \\[2pt]
$\iff$   & minimise the total surprisal $-\ell(\theta)$ \\[2pt]
$\iff$   & minimise $\bar{L}(\theta)$.
\end{tabular}
\end{center}
The first step holds because $\log$ is increasing. The
second is just a change of sign. The third holds because
dividing by $N$ does not change which $\theta$ is best.
This method of choosing $\theta$ is called
\emph{maximum likelihood estimation}.

\textbf{Example.} Take a coin model with
$\symbf{q} = (0.6,\, 0.4)$ for (Heads, Tails), and suppose
we observe Heads, Heads, Tails. Then
\begin{align*}
  \text{likelihood} &= 0.6 \times 0.6 \times 0.4 = 0.144, \\
  \text{log-likelihood} &= \log 0.144 \approx -1.938, \\
  \text{total surprisal} &= 0.511 + 0.511 + 0.916 \approx 1.938, \\
  \text{average surprisal} &\approx 1.938 / 3 \approx 0.646 \text{ nats}.
\end{align*}
The same number appears in every line, only with a change
of sign or a division by $N = 3$.

\paragraph{Who is who.}
Table~\ref{tab:names} puts the five names side by side.

\begin{table}[ht]
\centering
\begin{tabular}{llll}
  \toprule
  Name & What it is about & Needs observed data? & In training \\
  \midrule
  Surprisal & one outcome & no & a score \\
  Entropy & one distribution & no & a property \\
  Likelihood & the model's $\theta$ & yes & make large \\
  Log-likelihood & the model's $\theta$ & yes & make large \\
  Cross-entropy loss & the model's $\theta$ & yes & make small \\
  \bottomrule
\end{tabular}
\caption{Five names and how they relate. Negative
log-likelihood is total surprisal, and dividing it by $N$
gives the average cross-entropy loss.}
\label{tab:names}
\end{table}

\paragraph{Two points to keep in mind.}
\begin{enumerate}
  \item The likelihood is a function of $\theta$, with the
  data held fixed. It is \emph{not} a probability
  distribution over $\theta$: its values over different
  $\theta$ need not add up to $1$.
  \item Entropy $H(\symbf{q})$ belongs to a single
  distribution, and we can compute it without seeing any
  data. Likelihood, log-likelihood, and the cross-entropy
  loss always involve observed outcomes: they ask how well
  the model explains what actually happened.
\end{enumerate}

\section{From logits to $\symbf{q}$: the softmax}

A network does not output probabilities directly.
It outputs real numbers $z_1, \dots, z_K$,
called \emph{logits}. The softmax converts them
into a distribution:
\begin{equation}
  q_k = \frac{e^{z_k}}{\sum_{j} e^{z_j}} .
  \label{eq:softmax}
\end{equation}
Each $q_k$ is positive and they add up to $1$,
so $\symbf{q}$ is a valid distribution.

Putting this in the loss for true class $c$:
\begin{equation}
  L = -\log q_c = -z_c + \log \sum_{j} e^{z_j} .
  \label{eq:loss}
\end{equation}
\paragraph{In simple words:} the loss is small when
the correct logit $z_c$ is large compared with the other
logits. The log-sum-exp is then close to $z_c$, and the
two terms in Eq.~\eqref{eq:loss} nearly cancel. Raising
$z_c$ also raises the log-sum-exp, but by less (by $q_c$
per unit), so the loss still falls.

\paragraph{Note.} Deep learning libraries
compute $L$ directly from the logits using this
last form (with the log-sum-exp trick), instead
of first computing $q_k$ and then taking its log.
This avoids overflow in $e^{z_j}$ and avoids $\log 0$.

\paragraph{Gradient.} Differentiating with respect to
a logit gives
\begin{equation}
  \frac{\partial L}{\partial z_k} = q_k - p_k .
  \label{eq:grad}
\end{equation}
This is the model's probability minus the label.
Gradient descent moves each logit against its gradient.
For the true class the gradient is $q_c - 1 \le 0$, so the
logit is pushed up. For the wrong classes it
is $q_k \ge 0$, so those logits are pushed down.

\section{Summary}

\begin{center}
\begin{tabular}{lll}
  \toprule
  Eq. & Quantity & Meaning \\
  \midrule
  \eqref{eq:surprisal} & $-\log q_k$ & surprisal of outcome $k$
  under $\symbf{q}$ \\
  \eqref{eq:entropy} & $H(\symbf{q}) = -\sum_k q_k \log q_k$ & average surprisal
  under $\symbf{q}$ \\
  \eqref{eq:xent} & $H(\symbf{p}, \symbf{q}) = -\sum_k p_k \log q_k$ &
    average surprisal under $\symbf{q}$, data from $\symbf{p}$ \\
  \eqref{eq:onehot} & $-\log q_c$ & cross-entropy loss for one-hot label $c$ \\
  \eqref{eq:loglik} & $\ell(\theta) = \sum_i \log q_\theta(c_i \mid x_i)$ &
    log-likelihood of the data set \\
  \eqref{eq:meanloss} & $\bar{L}(\theta) = -\ell(\theta)/N$ &
    average cross-entropy loss \\
  \bottomrule
\end{tabular}
\end{center}

\end{document}